\chapter{状态测量、滤波与反馈时序}
公开数据采用TUM-VI\cite{tumvi}；滤波延迟的定义参照SciPy文档\cite{scipy}。


把虚拟传感器、公开实测数据和自有设备采集分别记录。采样、发布、控制与物理积分的频率不同，滤波平滑也可能增加数据滞后。

\section{反馈频率与虚拟传感器链路}

\label{stage:s05}

\subsection{原任务覆盖与成果来源}\label{s05-detail-1}

\FloatBarrier\begingroup\small\setlength{\tabcolsep}{2.5pt}
\begin{longtable}[]{@{}
  >{\raggedright\arraybackslash}p{(\linewidth - 4\tabcolsep) * \real{0.2451}}
  >{\raggedright\arraybackslash}p{(\linewidth - 4\tabcolsep) * \real{0.3859}}
  >{\raggedright\arraybackslash}p{(\linewidth - 4\tabcolsep) * \real{0.3690}}@{}}
\caption{原任务覆盖与成果来源：设置与结果}\label{tab:s05-1}\tabularnewline
\toprule\noalign{}
\begin{minipage}[b]{\linewidth}\raggedright
月度计划 Week 2 内容
\end{minipage} & \begin{minipage}[b]{\linewidth}\raggedright
该阶段完成内容
\end{minipage} & \begin{minipage}[b]{\linewidth}\raggedright
数据来源与边界
\end{minipage} \\
\midrule\noalign{}
\endfirsthead
\toprule\noalign{}
\begin{minipage}[b]{\linewidth}\raggedright
月度计划 Week 2 内容
\end{minipage} & \begin{minipage}[b]{\linewidth}\raggedright
该阶段完成内容
\end{minipage} & \begin{minipage}[b]{\linewidth}\raggedright
数据来源与边界
\end{minipage} \\
\midrule\noalign{}
\endhead
\bottomrule\noalign{}
\endlastfoot
读取并发布多轴状态 & 六轴 ROS2 Imu 原始/滤波节点与独立订阅验证；厂家 RAW 串口只读接口 & 发布验证使用明确标记的离线回放；真实串口分支尚未运行 \\
静止/运动原始采集 & Panda link7 上官方虚拟 IMU 的实际物理仿真记录 & 仿真读数，不是实体板卡实测 \\
滤波与平滑分析 & 5点移动平均、EMA、二阶 Butterworth；误差及滞后对照 & 在虚拟IMU基础上显式注入已知高斯噪声，参数不是实测估计 \\
多种反馈频率稳定性/延迟 & 25、50、100Hz各4次正反向平衡实验 & 真实ROS Bridge与两个外部ROS节点；物理步长固定10ms \\
频率影响比较 & 主机频率、间隔P95、消息完整性、链路延迟、跟踪误差 & 延迟定义限定链路，不声称完整硬件端到端响应 \\
\end{longtable}\endgroup

主实验 \path|runs/run002_balanced|；先导 \texttt{run001} 保留但不作为主要频率精度结论。先导9阶段存在运动方向分布不平衡，后续重跑12阶段，每个频率正/反各两次，减少该混杂因素。重复次数仍很少，同一物理场景，未做跨机器或负载条件统计。

\subsection{仿真及实际 ROS2 链路}\label{s05-detail-2}

复用 Week1 的 Panda 资产、五次轨迹和 Isaac 独立环境，新增专用话题，不覆盖原工作空间。

\begin{enumerate}
\def\labelenumi{\arabic{enumi}.}
\tightlist
\item
  Isaac Panda PhysX 100Hz
\item
  Bridge原始关节反馈
\item
  独立ROS中继 25/50/100Hz
\item
  独立ROS轨迹控制与记录节点
\item
  Bridge关节位置命令
\item
  link7官方虚拟IMU 100Hz
\item
  25Hz抽样及明确噪声注入
\item
  六轴因果滤波对照
\end{enumerate}

仿真10ms固定步长；中继按仿真时戳抽取反馈，控制节点在每个已接收反馈时更新位置目标。因此比较的是\textbf{外环状态反馈/命令更新频率}，不是改变 PhysX 内部驱动控制频率。状态及命令使用 \path|/panda_week2/|，ROS domain51、localhost、FastDDS。

每阶段2秒保持、5秒五次轨迹、2秒保持，共12阶段约108秒。前三关节端点同Week1，正反方向交替，其余关节目标保持。静止区指令不变，不假设传感器完全无噪声。原始反馈接收共10798条，相邻仿真时戳最大间隔10ms，未发现超过15ms的缺口。中继发送到控制订阅的序号无缺失/重复。

IMU挂载 \path|/World/Panda/panda_link7/Week2IMU|，局部姿态为单位四元数，读取三轴加速度、角速度和wxyz四元数；内部滤波宽度均1。加速度读取包含重力，静态不同姿态的重力投影会改变，不能当成平移加速度。传感器频率不能高于物理频率，故该阶段上限100Hz。

\subsection{反馈频率、误差与响应结果}\label{s05-detail-3}

表内为每频率4阶段指标的算术平均；``间隔P95平均''不是合并全部样本后的P95。关节RMSE只统计每阶段2--7秒运动窗口，比较同一时戳参考与实际位置，无滞后补偿。

\FloatBarrier\begingroup\footnotesize\setlength{\tabcolsep}{2pt}
\begin{longtable}[]{@{}
  >{\raggedright\arraybackslash}p{(\linewidth - 16\tabcolsep) * \real{0.1060}}
  >{\raggedleft\arraybackslash}p{(\linewidth - 16\tabcolsep) * \real{0.1060}}
  >{\raggedleft\arraybackslash}p{(\linewidth - 16\tabcolsep) * \real{0.1060}}
  >{\raggedleft\arraybackslash}p{(\linewidth - 16\tabcolsep) * \real{0.1161}}
  >{\raggedleft\arraybackslash}p{(\linewidth - 16\tabcolsep) * \real{0.1110}}
  >{\raggedleft\arraybackslash}p{(\linewidth - 16\tabcolsep) * \real{0.1163}}
  >{\raggedleft\arraybackslash}p{(\linewidth - 16\tabcolsep) * \real{0.1163}}
  >{\raggedleft\arraybackslash}p{(\linewidth - 16\tabcolsep) * \real{0.1163}}
  >{\raggedleft\arraybackslash}p{(\linewidth - 16\tabcolsep) * \real{0.1060}}@{}}
\caption{反馈频率、误差与响应结果：设置与结果}\label{tab:s05-2}\tabularnewline
\toprule\noalign{}
\begin{minipage}[b]{\linewidth}\raggedright
请求Hz
\end{minipage} & \begin{minipage}[b]{\linewidth}\raggedleft
次数
\end{minipage} & \begin{minipage}[b]{\linewidth}\raggedleft
实测主机Hz
\end{minipage} & \begin{minipage}[b]{\linewidth}\raggedleft
接收间隔P95平均 ms
\end{minipage} & \begin{minipage}[b]{\linewidth}\raggedleft
中继链路P95平均 ms
\end{minipage} & \begin{minipage}[b]{\linewidth}\raggedleft
J1 RMSE rad
\end{minipage} & \begin{minipage}[b]{\linewidth}\raggedleft
J2 RMSE rad
\end{minipage} & \begin{minipage}[b]{\linewidth}\raggedleft
J3 RMSE rad
\end{minipage} & \begin{minipage}[b]{\linewidth}\raggedleft
J1拟合滞后 ms
\end{minipage} \\
\midrule\noalign{}
\endfirsthead
\toprule\noalign{}
\begin{minipage}[b]{\linewidth}\raggedright
请求Hz
\end{minipage} & \begin{minipage}[b]{\linewidth}\raggedleft
次数
\end{minipage} & \begin{minipage}[b]{\linewidth}\raggedleft
实测主机Hz
\end{minipage} & \begin{minipage}[b]{\linewidth}\raggedleft
接收间隔P95平均 ms
\end{minipage} & \begin{minipage}[b]{\linewidth}\raggedleft
中继链路P95平均 ms
\end{minipage} & \begin{minipage}[b]{\linewidth}\raggedleft
J1 RMSE rad
\end{minipage} & \begin{minipage}[b]{\linewidth}\raggedleft
J2 RMSE rad
\end{minipage} & \begin{minipage}[b]{\linewidth}\raggedleft
J3 RMSE rad
\end{minipage} & \begin{minipage}[b]{\linewidth}\raggedleft
J1拟合滞后 ms
\end{minipage} \\
\midrule\noalign{}
\endhead
\bottomrule\noalign{}
\endlastfoot
25 & 4 & 24.891 & 40.964 & 0.259 & 0.021783 & 0.015586 & 0.016388 & 227.5 \\
50 & 4 & 49.639 & 21.129 & 0.214 & 0.020832 & 0.014911 & 0.015675 & 217.5 \\
100 & 4 & 99.270 & 10.891 & 0.166 & 0.020370 & 0.014582 & 0.015328 & 215.0 \\
\end{longtable}\endgroup

\stagefigure{figures/ch03/s05_frequency_comparison_d1863982.png}{反馈频率对照}{fig:s05-detail-1}

\begin{itemize}
\tightlist
\item
  中继链路延迟：同一机器 monotonic 时钟下，中继 publish 调用前时刻到控制节点接收时刻，包括DDS与调度，不包含 Isaac\(\rightarrow\)中继传输，不能写成整条机器人通信延迟。
\item
  拟合滞后：对参考轨迹搜索0--400ms、5ms步长的时间移位，求实际关节位置最接近的移位；是驱动动态与更新链路共同形成的描述量，不是纯网络延迟。表中的原始RMSE仍未移位。
\item
  100Hz的主机频率约99.27Hz，存在调度抖动，不能声称硬实时。提高外环频率后J1平均RMSE略降，但幅度有限，默认物理驱动仍产生明显跟踪滞后；没有证明仅提高频率可以消除误差。
\item
  改善仅针对同一仿真场景、默认驱动与本条轨迹。未引入力传感、导纳、姿态融合或真实机器人控制。
\end{itemize}


\paragraph{阶段结论}25/50/100Hz比较改变外环状态与命令更新，不改变100Hz物理驱动；频率提高没有消除默认驱动滞后。

\stagecode{s05}

\FloatBarrier

\section{因果滤波与ROS2消息验证}

\label{stage:s06}

\subsection{虚拟 IMU 因果滤波与噪声/滞后权衡}\label{s06-detail-1}

在实际100Hz虚拟IMU读数上每4帧抽一帧，得到25Hz数据，接近既有真实板卡历史输出频率。主实验抽样数2782；其中静止/运动评估区间在指标JSON中明确保存。加入独立零均值高斯噪声：加速度每轴\(\sigma\)=0.15m/s²，角速度每轴\(\sigma\)=0.02rad/s，固定seed20261013。噪声模型是人为设定，不是实体IMU标定结果。

只用当前与历史样本：

\begin{itemize}
\tightlist
\item
  移动平均：最近5个样本，完整窗口的线性相位延迟约80ms。
\item
  EMA：截止频率3Hz，\texttt{α=1-exp(-2πfcΔt)}；可用实际采样间隔更新，无前视。
\item
  二阶Butterworth：3Hz截止、25Hz设计采样率、SOS形式逐样本运行；没有使用非因果filtfilt。采样率变化时应重新设计系数。
\end{itemize}

误差为输出减去\textbf{同一时刻的干净虚拟IMU读数}。表中多轴RMSE汇总三个轴，降噪比例相对加入噪声的raw误差；没有把滤波后的曲线平滑程度当成精度。动态误差未做时间对齐，保留滤波延迟代价。

\FloatBarrier\begingroup\small\setlength{\tabcolsep}{2.5pt}
\begin{longtable}[]{@{}
  >{\raggedright\arraybackslash}p{(\linewidth - 12\tabcolsep) * \real{0.1326}}
  >{\raggedleft\arraybackslash}p{(\linewidth - 12\tabcolsep) * \real{0.1563}}
  >{\raggedleft\arraybackslash}p{(\linewidth - 12\tabcolsep) * \real{0.1326}}
  >{\raggedleft\arraybackslash}p{(\linewidth - 12\tabcolsep) * \real{0.1566}}
  >{\raggedleft\arraybackslash}p{(\linewidth - 12\tabcolsep) * \real{0.1326}}
  >{\raggedleft\arraybackslash}p{(\linewidth - 12\tabcolsep) * \real{0.1566}}
  >{\raggedleft\arraybackslash}p{(\linewidth - 12\tabcolsep) * \real{0.1326}}@{}}
\caption{虚拟 IMU 因果滤波与噪声/滞后权衡：设置与结果}\label{tab:s06-1}\tabularnewline
\toprule\noalign{}
\begin{minipage}[b]{\linewidth}\raggedright
算法
\end{minipage} & \begin{minipage}[b]{\linewidth}\raggedleft
静止加速度误差RMSE m/s²
\end{minipage} & \begin{minipage}[b]{\linewidth}\raggedleft
减少
\end{minipage} & \begin{minipage}[b]{\linewidth}\raggedleft
静止角速度误差RMSE rad/s
\end{minipage} & \begin{minipage}[b]{\linewidth}\raggedleft
减少
\end{minipage} & \begin{minipage}[b]{\linewidth}\raggedleft
运动角速度误差RMSE rad/s
\end{minipage} & \begin{minipage}[b]{\linewidth}\raggedleft
角速度最佳移位 ms
\end{minipage} \\
\midrule\noalign{}
\endfirsthead
\toprule\noalign{}
\begin{minipage}[b]{\linewidth}\raggedright
算法
\end{minipage} & \begin{minipage}[b]{\linewidth}\raggedleft
静止加速度误差RMSE m/s²
\end{minipage} & \begin{minipage}[b]{\linewidth}\raggedleft
减少
\end{minipage} & \begin{minipage}[b]{\linewidth}\raggedleft
静止角速度误差RMSE rad/s
\end{minipage} & \begin{minipage}[b]{\linewidth}\raggedleft
减少
\end{minipage} & \begin{minipage}[b]{\linewidth}\raggedleft
运动角速度误差RMSE rad/s
\end{minipage} & \begin{minipage}[b]{\linewidth}\raggedleft
角速度最佳移位 ms
\end{minipage} \\
\midrule\noalign{}
\endhead
\bottomrule\noalign{}
\endlastfoot
raw & 0.14903 & 0.0\% & 0.01983 & 0.0\% & 0.02002 & 0 \\
移动平均 & 0.06335 & 57.5\% & 0.00915 & 53.9\% & 0.00930 & 80 \\
指数平均 & 0.08870 & 40.5\% & 0.01203 & 39.3\% & 0.01197 & 40 \\
巴特沃斯 & 0.07292 & 51.1\% & 0.01028 & 48.1\% & 0.01027 & 80 \\
\end{longtable}\endgroup

\stagefigure{figures/ch03/s06_imu_filter_comparison_e8e8c507.png}{静止与运动滤波对照}{fig:s06-detail-1}

移动平均在此噪声和慢速轨迹条件下降噪较强，但引入约80ms滞后；EMA降噪较弱，拟合移位约40ms。移位以40ms采样网格计算，只用于描述，不抵消表中误差。常值与阶跃检查验证没有运动发生前的提前响应。因采样分辨率较粗，EMA首次变化即超过阶跃50\%并不意味着零相位延迟。

这轮优先选EMA作为演示发布默认值，保留其余算法对照；未给出所有运动速度下的最优参数，也未证明真实IMU同样改善。没有修改既有HTML姿态显示或板卡安装轴，没有对四元数做不合适的逐分量平均，当前六轴滤波与姿态融合分开。

\subsection{实际 ROS2 IMU 原始/滤波发布验证}\label{s06-detail-2}

另用先导实验带噪CSV做10秒离线回放，独立发布进程在线计算EMA，独立订阅进程记录 sensor\_msgs/Imu。CSV没有四元数字段，显式设置orientation不可用；没有假装获得融合姿态。

\FloatBarrier\begingroup\small\setlength{\tabcolsep}{2.5pt}
\begin{longtable}[]{@{}
  >{\raggedright\arraybackslash}p{(\linewidth - 8\tabcolsep) * \real{0.2409}}
  >{\raggedleft\arraybackslash}p{(\linewidth - 8\tabcolsep) * \real{0.1878}}
  >{\raggedleft\arraybackslash}p{(\linewidth - 8\tabcolsep) * \real{0.1859}}
  >{\raggedleft\arraybackslash}p{(\linewidth - 8\tabcolsep) * \real{0.1859}}
  >{\raggedleft\arraybackslash}p{(\linewidth - 8\tabcolsep) * \real{0.1994}}@{}}
\caption{实际 ROS2 IMU 原始/滤波发布验证：设置与结果}\label{tab:s06-2}\tabularnewline
\toprule\noalign{}
\begin{minipage}[b]{\linewidth}\raggedright
话题
\end{minipage} & \begin{minipage}[b]{\linewidth}\raggedleft
发送/接收
\end{minipage} & \begin{minipage}[b]{\linewidth}\raggedleft
缺失/重复
\end{minipage} & \begin{minipage}[b]{\linewidth}\raggedleft
主机Hz
\end{minipage} & \begin{minipage}[b]{\linewidth}\raggedleft
发布至接收P95 ms
\end{minipage} \\
\midrule\noalign{}
\endfirsthead
\toprule\noalign{}
\begin{minipage}[b]{\linewidth}\raggedright
话题
\end{minipage} & \begin{minipage}[b]{\linewidth}\raggedleft
发送/接收
\end{minipage} & \begin{minipage}[b]{\linewidth}\raggedleft
缺失/重复
\end{minipage} & \begin{minipage}[b]{\linewidth}\raggedleft
主机Hz
\end{minipage} & \begin{minipage}[b]{\linewidth}\raggedleft
发布至接收P95 ms
\end{minipage} \\
\midrule\noalign{}
\endhead
\bottomrule\noalign{}
\endlastfoot
\path|/imu_week2/raw| & 249/249 & 0/0 & 25.000 & 0.864 \\
\path|/imu_week2/filtered| & 249/249 & 0/0 & 25.000 & 0.969 \\
\end{longtable}\endgroup

原始/滤波消息采用固定坐标系名imu\_sensor，以相同ROS时间戳关联独立收发记录；六轴值与发布端记录最大差异0，两个进程正常退出。延迟起点为样本处理开始时刻，包含滤波计算、消息构建、DDS与调度；不包含CSV历史生成、原始传感器采样或串口延迟。此验证证明ROS接口能发送/接收，并非实体传感器联调。

真实串口分支已实现：使用原厂家驱动、115200接收有效RAW包、g\(\rightarrow\)m/s²、rad/s保留、可附上已收到的板卡四元数；不写入校准、输出频率或Flash。当前设备未连接，只完成代码检查，该分支尚待真实运行验证。


\paragraph{阶段结论}移动平均降低设定噪声但带来约80毫秒滞后；消息回放证明接口连通，与实体传感器联调分开。

\stagecode{s06}

\FloatBarrier

\section{公开真实IMU与噪声分析}

\label{stage:s07}

\subsection{公开真实IMU与六轴滤波}\label{s07-detail-1}

先比较了EuRoC和TUM-VI两类公开视觉惯性数据。EuRoC提供无人机上的200Hz六轴IMU，可用于后续跨设备对照；该阶段选择TUM-VI，是因为它另有独立静止记录，并有手持设备运动序列，更方便分别验证静止滤波和运动响应。EuRoC该阶段未下载或评测。

选择TUM-VI是因为它提供三轴加速度和角速度、原始时间戳，以及独立静止记录。数据由原作者采集，许可CC BY4.0。静止数据从原始NPY取前120001行，约602秒；滤波统计取其中60--90秒。运动段来自room1序列的中间固定30秒，未按结果挑选波形。通过官方服务器HTTP范围请求提取所需IMU记录，没有下载整套相机图像。

\FloatBarrier\begingroup\small\setlength{\tabcolsep}{2.5pt}
\begin{longtable}[]{@{}
  >{\raggedright\arraybackslash}p{(\linewidth - 6\tabcolsep) * \real{0.2359}}
  >{\raggedright\arraybackslash}p{(\linewidth - 6\tabcolsep) * \real{0.2359}}
  >{\raggedright\arraybackslash}p{(\linewidth - 6\tabcolsep) * \real{0.2437}}
  >{\raggedright\arraybackslash}p{(\linewidth - 6\tabcolsep) * \real{0.2845}}@{}}
\caption{公开真实IMU与六轴滤波：设置与结果}\label{tab:s07-1}\tabularnewline
\toprule\noalign{}
\begin{minipage}[b]{\linewidth}\raggedright
数据段
\end{minipage} & \begin{minipage}[b]{\linewidth}\raggedright
样本数
\end{minipage} & \begin{minipage}[b]{\linewidth}\raggedright
实际覆盖 / s
\end{minipage} & \begin{minipage}[b]{\linewidth}\raggedright
时间戳统计频率 / Hz
\end{minipage} \\
\midrule\noalign{}
\endfirsthead
\toprule\noalign{}
\begin{minipage}[b]{\linewidth}\raggedright
数据段
\end{minipage} & \begin{minipage}[b]{\linewidth}\raggedright
样本数
\end{minipage} & \begin{minipage}[b]{\linewidth}\raggedright
实际覆盖 / s
\end{minipage} & \begin{minipage}[b]{\linewidth}\raggedright
时间戳统计频率 / Hz
\end{minipage} \\
\midrule\noalign{}
\endhead
\bottomrule\noalign{}
\endlastfoot
静止 & 5981 & 29.994 & 199.374 \\
运动 & 5981 & 29.995 & 199.368 \\
\end{longtable}\endgroup

六通道分别使用当前及过去四点的因果平均，不使用未来数据。前四点窗口尚未完整，均值和标准差比较统一排除这四点。加速度保持m/s²、角速度保持rad/s，不把静止加速度中的重力当成零偏清除。

静止段统计如下：

\FloatBarrier\begingroup\small\setlength{\tabcolsep}{2.5pt}
\begin{longtable}[]{@{}
  >{\raggedright\arraybackslash}p{(\linewidth - 10\tabcolsep) * \real{0.1598}}
  >{\raggedright\arraybackslash}p{(\linewidth - 10\tabcolsep) * \real{0.1598}}
  >{\raggedright\arraybackslash}p{(\linewidth - 10\tabcolsep) * \real{0.1710}}
  >{\raggedright\arraybackslash}p{(\linewidth - 10\tabcolsep) * \real{0.1710}}
  >{\raggedright\arraybackslash}p{(\linewidth - 10\tabcolsep) * \real{0.1692}}
  >{\raggedright\arraybackslash}p{(\linewidth - 10\tabcolsep) * \real{0.1692}}@{}}
\caption{公开真实IMU与六轴滤波：设置与结果}\label{tab:s07-2}\tabularnewline
\toprule\noalign{}
\begin{minipage}[b]{\linewidth}\raggedright
通道
\end{minipage} & \begin{minipage}[b]{\linewidth}\raggedright
单位
\end{minipage} & \begin{minipage}[b]{\linewidth}\raggedright
原始均值
\end{minipage} & \begin{minipage}[b]{\linewidth}\raggedright
滤波均值
\end{minipage} & \begin{minipage}[b]{\linewidth}\raggedright
原始标准差
\end{minipage} & \begin{minipage}[b]{\linewidth}\raggedright
滤波标准差
\end{minipage} \\
\midrule\noalign{}
\endfirsthead
\toprule\noalign{}
\begin{minipage}[b]{\linewidth}\raggedright
通道
\end{minipage} & \begin{minipage}[b]{\linewidth}\raggedright
单位
\end{minipage} & \begin{minipage}[b]{\linewidth}\raggedright
原始均值
\end{minipage} & \begin{minipage}[b]{\linewidth}\raggedright
滤波均值
\end{minipage} & \begin{minipage}[b]{\linewidth}\raggedright
原始标准差
\end{minipage} & \begin{minipage}[b]{\linewidth}\raggedright
滤波标准差
\end{minipage} \\
\midrule\noalign{}
\endhead
\bottomrule\noalign{}
\endlastfoot
ax & m/s² & -1.416228 & -1.416217 & 0.017933 & 0.008507 \\
ay & m/s² & -0.396287 & -0.396288 & 0.013538 & 0.006626 \\
az & m/s² & 10.487650 & 10.487644 & 0.020199 & 0.009757 \\
gx & rad/s & 0.029597 & 0.029597 & 0.001105 & 0.000554 \\
gy & rad/s & 0.006676 & 0.006676 & 0.000947 & 0.000481 \\
gz & rad/s & 0.016164 & 0.016164 & 0.000988 & 0.000493 \\
\end{longtable}\endgroup

六通道标准差降低约49.2\%--52.6\%。5点平均的中心相对最新采样落后两个采样周期，当前原始数据约为10.03ms；若输入只有25Hz，同样5点窗口约为80ms。平滑程度不能脱离输入采样率比较。

\stagefigure{figures/ch03/s07_static_six_axes_126abe0a.png}{静止六轴滤波}{fig:s07-detail-1}

\stagefigure{figures/ch03/s07_motion_six_axes_642cb4e5.png}{运动六轴滤波}{fig:s07-detail-2}

静止均值也不作为已标定零偏；这里没有依据当前姿态的真值去分离重力、零偏和比例误差。

运动段同时保存均值和标准差表，但运动信号本身会变化，标准差降低不能直接解释为噪声或测量误差下降。没有同步运动真值，当前不报告滤波后的角度或位移精度。

相关代码在 \path|scripts/download_public_imu.py|、\path|scripts/download_motion.py|、\path|scripts/analyze_imu.py|、\path|scripts/filter_csv.py|；运动段统计在 \path|runs/motion_filter_statistics.csv|。

\subsection{静止频谱与Allan曲线}\label{s07-detail-2}

在约602秒静止数据上补充频谱密度和重叠Allan偏差。前者观察不同频率上的波动分布；后者观察不同平均时间下的变化，平均时间最长约30秒。数据同时保留温度，当前范围约31.39--33.00℃。

\stagefigure{figures/ch03/s07_static_psd_21a3cffd.png}{静止数据频谱}{fig:s07-detail-3}

\stagefigure{figures/ch03/s07_static_allan_e2029130.png}{静止数据Allan曲线}{fig:s07-detail-4}

这些图用于探索窗口长度和长时间变化，不等于完成传感器标定。当前截取长度有限，没有拟合或报告硬件噪声密度、零偏不稳定性等标定参数，也不将公开传感器参数用于自有板卡。

相关代码在 \path|scripts/analyze_imu.py|。


\paragraph{阶段结论}TUM-VI六轴静止标准差下降约49.2\%至52.6\%；缺运动真值，不能把运动波形变平解释成测量精度提高。

\stagecode{s07}

\FloatBarrier

\section{发布频率、数据年龄与Jetson实体采集}

\label{stage:s08}

\subsection{50/100/200Hz反馈测试}\label{s08-detail-1}

三个独立ROS进程分别负责数据回放、5点平均与定频发布、状态接收。原始及滤波消息采用 \path|sensor_msgs/msg/Imu|，话题为 \path|/imu/data_raw|、\path|/imu/data|。额外消息记录输入序号、发布序号及同一主机上的单调时钟，方便区分重复测量、主动降采样和真正的消息缺失。

输入设为名义200Hz或25Hz，输出设为50、100、200Hz，每组约30秒。低频输入由公开运动记录每8点取1点，不是自有IMU的实测采样。回放按软件定时器发送，原始记录实际约199.37Hz，不能把回放发布率当成原传感器频率。统计取稳定运行的第1至29秒，排除启动和停止过渡。

\FloatBarrier\begingroup\small\setlength{\tabcolsep}{2.5pt}
\begin{longtable}[]{@{}
  >{\raggedright\arraybackslash}p{(\linewidth - 12\tabcolsep) * \real{0.1429}}
  >{\raggedright\arraybackslash}p{(\linewidth - 12\tabcolsep) * \real{0.1429}}
  >{\raggedright\arraybackslash}p{(\linewidth - 12\tabcolsep) * \real{0.1429}}
  >{\raggedright\arraybackslash}p{(\linewidth - 12\tabcolsep) * \real{0.1429}}
  >{\raggedright\arraybackslash}p{(\linewidth - 12\tabcolsep) * \real{0.1429}}
  >{\raggedright\arraybackslash}p{(\linewidth - 12\tabcolsep) * \real{0.1429}}
  >{\raggedright\arraybackslash}p{(\linewidth - 12\tabcolsep) * \real{0.1429}}@{}}
\caption{50/100/200Hz反馈测试：设置与结果}\label{tab:s08-1}\tabularnewline
\toprule\noalign{}
\begin{minipage}[b]{\linewidth}\raggedright
输入 / Hz
\end{minipage} & \begin{minipage}[b]{\linewidth}\raggedright
目标反馈 / Hz
\end{minipage} & \begin{minipage}[b]{\linewidth}\raggedright
实测反馈 / Hz
\end{minipage} & \begin{minipage}[b]{\linewidth}\raggedright
间隔标准差 / ms
\end{minipage} & \begin{minipage}[b]{\linewidth}\raggedright
数据年龄P95 / ms
\end{minipage} & \begin{minipage}[b]{\linewidth}\raggedright
重复测量 / \%
\end{minipage} & \begin{minipage}[b]{\linewidth}\raggedright
发送\(\rightarrow\)接收缺失
\end{minipage} \\
\midrule\noalign{}
\endfirsthead
\toprule\noalign{}
\begin{minipage}[b]{\linewidth}\raggedright
输入 / Hz
\end{minipage} & \begin{minipage}[b]{\linewidth}\raggedright
目标反馈 / Hz
\end{minipage} & \begin{minipage}[b]{\linewidth}\raggedright
实测反馈 / Hz
\end{minipage} & \begin{minipage}[b]{\linewidth}\raggedright
间隔标准差 / ms
\end{minipage} & \begin{minipage}[b]{\linewidth}\raggedright
数据年龄P95 / ms
\end{minipage} & \begin{minipage}[b]{\linewidth}\raggedright
重复测量 / \%
\end{minipage} & \begin{minipage}[b]{\linewidth}\raggedright
发送\(\rightarrow\)接收缺失
\end{minipage} \\
\midrule\noalign{}
\endhead
\bottomrule\noalign{}
\endlastfoot
200 & 50 & 50.00 & 0.196 & 6.065 & 0.0 & 0 \\
200 & 100 & 100.00 & 0.241 & 4.552 & 0.0 & 0 \\
200 & 200 & 200.00 & 0.087 & 1.485 & 0.0 & 0 \\
25 & 50 & 50.00 & 0.344 & 34.253 & 50.0 & 0 \\
25 & 100 & 100.00 & 0.349 & 35.710 & 75.0 & 0 \\
25 & 200 & 200.00 & 0.287 & 38.196 & 87.5 & 0 \\
\end{longtable}\endgroup

``数据年龄''从输入样本在回放进程被取出，到接收进程收到对应反馈为止；还未包含5点平均相对中心的滞后，更不包含真实传感器内部与串口延迟。单独的发布到接收耗时和P95保存在metrics.csv。接收节点只检查并接收状态，不构成执行器端的物理响应时间。

25Hz输入提高到200Hz发布后，大量消息重复使用同一测量；更高的发布率没有增加新的传感器信息。200Hz输入降为50Hz时，有些中间样本未作为反馈发出，这是主动取最新值，不能算作网络丢帧。当前结果只针对本机有线软件进程与所设QoS，不推广为无线或Jetson性能。

\stagefigure{figures/ch03/s08_feedback_rates_856eba1b.png}{反馈频率与数据新鲜度}{fig:s08-detail-1}

相关代码在 \path|scripts/feedback_node.py|、\path|scripts/run_feedback.py| 和 \path|scripts/analyze_feedback.py|；详细结果在 \path|runs/feedback001/metrics.csv|。

\subsection{Jetson实体IMU测试}\label{s08-detail-2}

该次把此前 PC 端流程迁移到 Jetson，完成实体 IMU 的静止、运动采集及三档 ROS 2 发布测试。没有启动机械臂，也没有修改传感器输出配置。

\subsubsection{环境与方法}\label{s08-detail-3}

Jetson Orin Nano，Ubuntu 22.04.5、L4T R36.4.7、ROS 2 Humble。使用系统 Python 3.10.12，不修改现有 an-project Conda 环境。串口为 \path|/dev/serial/by-id/usb-1a86_USB_Serial-if00-port0|，115200 波特率。

采集程序解析 RAW 六轴数据，将加速度转换为 m/s²，角速度为 rad/s；每个新样本更新一次 5 点因果均值。独立进程按目标频率发布 \path|/imu/data|，另一个进程接收。全部测时进程位于同一 Jetson，使用同机单调时钟。姿态未参与本实验，ROS 消息将 orientation\_covariance{[}0{]} 设为 -1。

\subsubsection{静止与运动采集}\label{s08-detail-4}

两段各采集约 30 秒、749 组。静止首末样本估计 25.131 Hz，运动 25.052 Hz。静止加速度模长约 9.84 m/s²。运动包含三个方向的转动；三轴角速度绝对峰值约 4.948、4.412、3.129 rad/s。在线滤波值与离线重算一致。静止 gx、gy 为零，可能受设备量化或内部处理影响，不能解读为无噪声。

\stagefigure{figures/ch03/s08_six_axes_39c424bd.png}{静止六轴}{fig:s08-detail-2}

\stagefigure{figures/ch03/s08_six_axes_b188f7b9.png}{运动六轴}{fig:s08-detail-3}

\subsubsection{发布频率测试}\label{s08-detail-5}

每档直接读取实体 IMU 约 30 秒，分析中间稳定区间，排除启动和停止尾段。

\FloatBarrier\begingroup\small\setlength{\tabcolsep}{2.5pt}
\begin{longtable}[]{@{}
  >{\raggedright\arraybackslash}p{(\linewidth - 12\tabcolsep) * \real{0.1395}}
  >{\raggedright\arraybackslash}p{(\linewidth - 12\tabcolsep) * \real{0.1395}}
  >{\raggedright\arraybackslash}p{(\linewidth - 12\tabcolsep) * \real{0.1395}}
  >{\raggedright\arraybackslash}p{(\linewidth - 12\tabcolsep) * \real{0.1395}}
  >{\raggedright\arraybackslash}p{(\linewidth - 12\tabcolsep) * \real{0.1395}}
  >{\raggedright\arraybackslash}p{(\linewidth - 12\tabcolsep) * \real{0.1399}}
  >{\raggedright\arraybackslash}p{(\linewidth - 12\tabcolsep) * \real{0.1626}}@{}}
\caption{发布频率测试：设置与结果}\label{tab:s08-2}\tabularnewline
\toprule\noalign{}
\begin{minipage}[b]{\linewidth}\raggedright
目标 Hz
\end{minipage} & \begin{minipage}[b]{\linewidth}\raggedright
实际 Hz
\end{minipage} & \begin{minipage}[b]{\linewidth}\raggedright
新测量 Hz
\end{minipage} & \begin{minipage}[b]{\linewidth}\raggedright
重复比例
\end{minipage} & \begin{minipage}[b]{\linewidth}\raggedright
间隔标准差 ms
\end{minipage} & \begin{minipage}[b]{\linewidth}\raggedright
消息传输 P95 ms
\end{minipage} & \begin{minipage}[b]{\linewidth}\raggedright
最新原始样本年龄 P95 ms
\end{minipage} \\
\midrule\noalign{}
\endfirsthead
\toprule\noalign{}
\begin{minipage}[b]{\linewidth}\raggedright
目标 Hz
\end{minipage} & \begin{minipage}[b]{\linewidth}\raggedright
实际 Hz
\end{minipage} & \begin{minipage}[b]{\linewidth}\raggedright
新测量 Hz
\end{minipage} & \begin{minipage}[b]{\linewidth}\raggedright
重复比例
\end{minipage} & \begin{minipage}[b]{\linewidth}\raggedright
间隔标准差 ms
\end{minipage} & \begin{minipage}[b]{\linewidth}\raggedright
消息传输 P95 ms
\end{minipage} & \begin{minipage}[b]{\linewidth}\raggedright
最新原始样本年龄 P95 ms
\end{minipage} \\
\midrule\noalign{}
\endhead
\bottomrule\noalign{}
\endlastfoot
50 & 50.000 & 24.967 & 50.04\% & 0.497 & 1.095 & 40.109 \\
100 & 99.999 & 24.966 & 75.03\% & 0.520 & 0.816 & 39.605 \\
200 & 200.001 & 24.968 & 87.52\% & 0.390 & 0.759 & 40.164 \\
\end{longtable}\endgroup

三档共发布并接收 10494 条反馈消息，元数据与 Imu 消息计数一致，发布至接收的序号未发现缺口。板卡协议没有在此解析设备样本序号，因此不能据此认定传感器或 USB 链路零丢样。

消息传输指 relay 发布前至 sink 回调接收的时间；样本年龄从主机解析 RAW 包开始计算。两者均不是设备采样到控制执行的完整时延，也不包含 5 点滤波的约 80 ms 群延迟。提高发布频率主要减少等待下一次发布的时间，不能补足新的传感器信息。

\subsubsection{文件与后续工作}\label{s08-detail-6}

相关代码在 \path|scripts/imu_serial_node.py|、\path|scripts/feedback_node.py|、\path|scripts/run_feedback.py|、\path|scripts/filter_csv.py|、\path|scripts/analyze_feedback.py|；启动环境由 \path|scripts/ros-python| 指定。相对路径均以本报告所在目录为根。

原始数据与分析位于 \path|runs/jetson_20261020/|：

\begin{Verbatim}[breaklines=true,breakanywhere=true,fontsize=\footnotesize]
hardware_static_20261020_01/   静止 CSV、原始串口字节、会话信息
hardware_motion_20261020_02/   运动 CSV、原始串口字节、会话信息
static_analysis_01/           六轴图、均值与标准差、滤波 CSV
motion_analysis_02/           六轴图、均值与标准差、滤波 CSV
hardware_feedback_20261020_01/
  source25_feedback50/        三进程日志、CSV、原始字节
  source25_feedback100/
  source25_feedback200/
  metrics.csv                发布频率、间隔、时延与重复比例
  feedback_rates.png
\end{Verbatim}

Jetson 同名工程目录保留原始记录；上述资料已回传工作站。此前暂停的 motion\_01 不用于该次分析。后续先讨论约 25 Hz 的设备输出与 80 ms 滤波延迟是否适合目标控制任务，再决定是否调整设备配置或滤波窗口；当前未改固件参数，也未运行机械臂闭环。


\paragraph{阶段结论}实体IMU约25Hz产生新测量，200Hz发布时约87.52\%重复；传输P95不足以代表包含滤波的端到端时延。

\stagecode{s08}

\FloatBarrier
